\section{第41届复赛第4题：高能电子探测与质子散射的不确定度估计}
\subsection{题干}
第（1）问以质子半径 $r=0.850\,\mathrm{fm}$ 为尺度，要求用电子完成200层断层扫描，确定电子的最小动能，并比较狭义相对论和极端相对论近似。第（2）问讨论质量为 $m_0$、电荷为 $Z_1e$ 的入射粒子与电荷为 $Z_2e$ 的固定靶之间的非相对论散射；要求写出总能量，并在两个质子中心距离的测量不确定度约等于该距离时，估计入射质子的速率。
\subsection{参考解答整理}
200层扫描对应空间分辨率 $\Delta x=2r/200=r/100$。由不确定关系与相对论能量动量关系，
\[\Delta x\,\Delta p_x\geq\frac{\hbar}{2},\qquad E_{k,\min}=\sqrt{\left(\frac{\hbar c}{2\Delta x}\right)^2+m_e^2c^4}-m_ec^2.\]
代入 $\Delta x=r/100$ 后，答案给出的最小电子动能为 \[\boxed{E_{k,\min}\approx11.6\,\mathrm{GeV}}.\]
极端相对论近似为 $E_{k,\min}\approx\hbar c/(2\Delta x)$，与狭义相对论结果的差别远小于1\%。答案第9式的焦耳单位指数与同式的 GeV 值不一致，故此处不转写该冲突的焦耳数，仅保留两种表达式共同支持的 GeV 结论。
\begin{figure}[htbp]\centering\begin{tikzpicture}[>=Stealth,scale=.88]
  \shade[ball color=red!45] (0,0) circle(1.25); \node at (0,0) {质子};
  \draw[<->] (-1.25,-1.55)--(1.25,-1.55) node[midway,below] {$2r=1.70\,\mathrm{fm}$};
  \foreach \x in {-1.0,-.5,0,.5,1.0} {\draw[densely dashed] (\x,-1.1)--(\x,1.1);}
  \draw[<->] (-1.0,1.58)--(-.5,1.58) node[midway,above] {$\Delta x=2r/200$};
  \draw[->,very thick,blue!70!black] (-4.0,0)--(-1.55,0) node[midway,above] {电子束};
  \draw[decorate,decoration={snake,amplitude=1.2pt,segment length=7pt},blue!65!black] (-3.8,.35)--(-1.55,.35);
  \node[align=center] at (3.25,.45) {$p_{\min}\geq\frac{\hbar}{2\Delta x}$\\[-2pt]$E_k\approx11.6\,\mathrm{GeV}$};
  \node[below,align=center] at (0,-2.2) {解析辅助图：200层断层扫描的空间分辨率与电子动量尺度};
\end{tikzpicture}
\caption{第41届复赛第4题解析辅助图：断层扫描分辨率与电子动量尺度。}\end{figure}

第（2）问中系统总能量为
\[E_t=\frac{p^2}{2m_0}+\frac{Z_1Z_2e^2}{4\pi\varepsilon_0r}.\]
令 $p=m_0v$，结合 $\Delta p\geq\hbar/(2\Delta r)$ 和能量守恒条件，得到
\[\frac{\Delta r}{r}=\sqrt{\frac{\varepsilon_0vh}{Z_1Z_2e^2}}.\]
对质子--质子情形 $Z_1=Z_2=1$，且 $\Delta r/r\simeq1$，答案给出
\[\boxed{v\approx0.044\times10^8\,\mathrm{m\,s^{-1}}\approx0.015c}.\]
\begin{figure}[htbp]\centering\begin{tikzpicture}[>=Stealth,scale=.9]
  \draw[->,thick] (0,0)--(6.3,0) node[right] {$r$};
  \draw[->,thick] (0,0)--(0,3.9) node[above] {$E$};
  \draw[very thick,blue!70!black,domain=.75:5.8,samples=80] plot (\x,{.38*\x+1.5/\x});
  \draw[densely dashed] (2,0)--(2,1.51); \fill (2,1.51) circle(1.4pt);
  \node[below] at (2,0) {$\Delta r/r\simeq1$};
  \node[above right,align=left] at (2.1,1.55) {不确定性动能项与\\库仑势能项的平衡};
  \node[right] at (4.4,3.25) {$E_t=\frac{p^2}{2m_0}+\frac{Z_1Z_2e^2}{4\pi\varepsilon_0r}$};
  \node[below,align=center] at (2.8,-.72) {解析辅助图：位置不确定度约束下的能量平衡};
\end{tikzpicture}
\caption{第41届复赛第4题解析辅助图：位置不确定度约束下的能量平衡。}\end{figure}
\subsection{转写依据}
题干依据 \texttt{CPHO41-Q04\_transcription\_v1.md}；答案公式与数值依据答案第9至11页及 \texttt{CPHO41-Q04\_formula\_check\_v1.csv}。原题没有独立题图；两幅图均为解析辅助图。极端相对论近似的焦耳值存在同页数值--单位冲突，已在本条目中保留为待核实边界，不据此扩展结论。
